\section{The Growth of Christian Hope}

Christian hope is a theological virtue: through it we desire eternal life and trust the promises of Christ, relying upon the help of the Holy Spirit. The Catechism presents it as God's answer to our aspiration for happiness, directing that aspiration toward his kingdom (CCC 1817--1821). Hope can therefore remain alive when we cannot predict a favorable earthly outcome. Its assurance rests upon the Lord's fidelity and power to save. We may be unable to see how a particular difficulty will end, yet still place our lives before the God who calls us to communion. The virtue enables a real movement toward him within that uncertainty.

Hope receives the demands of perseverance together with the grace by which we can answer them. The Catechism joins the expectation of heaven to Christ's merits, prayer, and the patience needed during trial. A difficult day need not become a verdict upon the whole pilgrimage. We can ask for the help required to remain faithful within it. Prayer draws expectation into the present: the kingdom we await becomes the good for which we ask, and the Father's will becomes the order to which we consent. The promised end sustains the next act of obedience without making the entire course of our future visible in advance.

Bernard's account of growing love helps us understand how this hope can become more generous. In \emph{On the Love of God} he begins with the human person who naturally attends to his own needs (VIII). Bodily weakness and dependence are facts of our condition. Hunger requires food; danger calls for protection; trouble brings us to seek help. Bernard does not expect a creature to begin as though it had no need to be sustained. He asks how the love that begins with ourselves can be drawn into a larger order. The command to love our neighbor as ourselves checks the impulse to make our own advantage the measure of every good.

Dependence can bring a person to seek God. The person who needs help can learn that he is neither his own creator nor his own sufficient keeper. He can turn to God because a need has become urgent and gradually recognize the goodness of the one to whom he turns. An imperfect beginning need not be despised. Many prayers first take shape through fear, illness, or the need of someone loved. What matters is whether we allow the encounter to deepen. We can receive an answer gratefully, acknowledge the giver, and remain open to his will beyond the particular relief we sought. Need becomes a place where the knowledge of God may begin to grow.

Bernard describes a further movement from loving God for the help he gives to loving him for his own goodness (IX). Frequent recourse to the divine helper brings a fuller acquaintance with the one who helps. His account is personal: the gift opens a relationship in which the giver becomes loved. A friendship among creatures can offer a limited analogy. We may first seek someone because he can teach us a skill, then come to value his wisdom and finally his person. In the life of grace the movement reaches an altogether greater object. God draws us to discover that his goodness deserves love beyond the occasions on which it has relieved an immediate need.

The distinction helps us examine our expectation of heaven. We rightly desire safety from sin, deliverance from death, and the joy of a life made whole. Yet these promises can remain loosely connected to the living God in our imagination. Bernard invites us to let the desire deepen until we love the promised communion because it is communion with him. The goodness revealed in Christ becomes the reason we want to be with Christ. Heaven is then increasingly desired as the fulfillment of a friendship already begun. We welcome the gifts through which that friendship is sustained, while learning to seek its giver in everything he allows us to receive.

Bernard's account also enlarges love of neighbor. When God is loved for himself, the people who belong to him enter more deeply into our concern. Charity is freed from the narrow circle of those whose companionship directly benefits us. We can seek another person's good because he is God's creature and is called to God's friendship. This gives hope a communal direction. We desire salvation with others and become willing to serve their journey. A household's hospitality can express that desire by receiving someone whose need brings no immediate advantage to its members. They serve a person called to God's friendship and begin to welcome his good as their own concern.

In the following chapter Bernard considers loving oneself for God's sake (X). His sequence helps us see why self-forgetfulness need not mean self-contempt. The creature's own good is brought within the goodness he loves in God. We may ask for our salvation wholeheartedly because God wills to give himself to us, and we may welcome the transformation through which we become capable of receiving that gift. Bernard's account of degrees supplies a direction for desire. We need not turn it into a method for assigning ourselves a spiritual rank. Our motives remain mixed, and God can lead us through many returns to elementary dependence. The useful question is whether love is becoming more truthful, grateful, and ready to consent to him.

Bernard has already considered the relation of love to reward (VII). A love directed toward something other than its beloved would be valuing the beloved chiefly as a means. The reward proper to loving God is God himself. This keeps the desire for eternal happiness within charity: the person who hopes for God hopes to receive the one whom he loves and to love him more perfectly. Our good is fulfilled in belonging to him. As charity grows, we come to desire his presence for the goodness we have learned to love, and the reward is the fuller possession of that same good.

When we return to this meditation, we may discover that our longing for heaven still depends heavily upon escape from a particular burden. Bernard's account allows us to recognize this mixed beginning and ask to know and love God more fully within it. We may also discover a new readiness to give thanks, or a freedom to serve without demanding an immediate return. These are gifts to receive gratefully, rather than measures of our whole spiritual condition. As need recurs, we can return to God with a clearer knowledge of his goodness and a greater readiness to receive his will.

